\section{Related Works}
\label{sec:rel}

\paragraph{Sharded blockchains.}
BlockchainDB, ShaPer, ByShard and GriDB are sharded blockchains with a similar security model to \sys.
They propose to deploy multiple consensus instances on different sets of nodes, each tolerating up to $f$ Byzantine faults.
Other sharded blockchains, such as Omniledger, RapidChain, and Dang et al., adopt a different security model where the systems tolerate a certain fraction of Byzantine nodes across the entire system, and the node shards are formed via random sampling at runtime.
As discussed in \cref{sec:review}, both categories of sharded blockchains suffer from performance degradation due to cross-shard transactions.
Basil~\cite{suri2021basil} discussed more shortcomings of sharded blockchains in detail.

\paragraph{BFT scalability.}
Many works have proposed to scale BFT consensus protocols, including SBFT, HotStuff, Tusk, Bullshark, and Autobahn.
Some of them also adopt sharding techniques, such as Elastico and OHIE.
While these works significantly improve the scalability of BFT consensus, the transaction execution is still sequential and fully replicated, becoming the bottleneck of the system.

\paragraph{Blockchain execution engines.}
LVMT and Block-STM improve the performance of blockchain execution engines by optimize IO efficiency and parallelizing execution, respectively.
They improve the performance of individual nodes and can be used together with \sys as scalably performant BFT systems.